\documentclass[10pt,a4paper]{amsart}
\usepackage[margin=19mm]{geometry}
\usepackage{amsmath,amssymb,mathtools}
\usepackage[scr=rsfs,cal=euler]{mathalfa}
\usepackage{xfrac}
\usepackage[unicode,hidelinks,bookmarks=false]{hyperref}
\newcommand{\ie}{i.e.\ }
\newcommand{\cf}{cf.\ }
\newcommand{\resp}{respectively\ }
\newcommand{\Subsection}{\subsection}
\providecommand{\nobreakdash}{\nobreak\hbox{-}}
\setlength{\emergencystretch}{2em}
\multlinegap0pt
\usepackage{bm}
\usepackage{bbm}
\usepackage{dsfont}
\usepackage{xspace}
\usepackage{cancel}
\usepackage{mathtools}
\usepackage{amsthm}
\makeatletter
\@ifundefined{thm}{\newtheorem{thm}{Thm}}{}
\@ifundefined{prop}{\newtheorem{prop}[thm]{Proposition}}{}
\makeatother
\DeclareMathOperator{\conv}{conv}
\newcommand{\bb}[1]{\mathbb{#1}}
\newcommand{\cl}[1]{\mathcal{#1}}
\title[Unique Games and Games Based on Groups]{Unique Games and Games Based on Groups}
\author{Rupert H. Levene, Vern I. Paulsen}
\date{Source: arXiv:2506.18644v1 (CC BY 4.0)}
\begin{document}
\maketitle

\noindent\emph{Source and licence.} Unique Games and Games Based on Groups. Version arXiv:2506.18644v1 (CC BY 4.0), distributed under the Creative Commons Attribution 4.0 International licence, as recorded in the arXiv licence metadata for this exact version and shown on the abstract page https://arxiv.org/abs/2506.18644v1. Full rights notice: licenses/arxiv-2506.18644-CC-BY-4.0.txt.\par\smallskip

\noindent\textit{Editorial scope.} Counted unit: the Prop whose source label is \texttt{None} in the article (source0.tex lines 665--667, source label \texttt{None}), delivered together with the article's own complete original proof of it (source0.tex lines 668--771, one \texttt{proof} environment of 104 lines). No mathematics is written, repaired or re-derived: statement and proof are the article's own text line for line. Required context retained uncounted: none. Every definition, display and label of those lines is retained unchanged. Omitted from this package and not redistributed: 1--664, 772--948. Nothing inside a retained excerpt is omitted or altered: every retained body is delivered exactly as the article's lines are written, so no omission receipt of any kind accompanies this package. Citations: the retained excerpts carry no citation command at all, so no bibliography entry of the article is transcribed and no reference is redistributed. Reference adaptation: none was needed and none was made; every reference inside the delivered unit resolves inside it, so no unresolved reference or citation is produced. Printed slips kept literal: no printed word, symbol, hypothesis or number of the counted statement or of its proof was altered. Layout and class: the article's own class is \texttt{amsart} with packages including fontenc, amsmath, amsfonts, amssymb, hyperref; the wrapper is the lane's pinned \texttt{amsart} editorial wrapper at 10pt on A4, which restates the theorem-like environments the retained text needs and adds the article's own macro definitions that the retained text needs (\texttt{\textbackslash bb}, \texttt{\textbackslash cl}, \texttt{\textbackslash conv}), with extra wrapper packages (bm, bbm, dsfont, xspace, cancel, mathtools). The delivered numbering is the unit's own. Assets: no retained span contains a figure environment, a graphics-inclusion command, a PDF-page inclusion, a table or a TikZ picture; no image file is copied into this package, the package declares no asset dependency and no image file is redistributed with it. Licence: version arXiv:2506.18644v1 of this article is distributed under the Creative Commons Attribution 4.0 International licence, as recorded in the arXiv licence metadata for this exact version and shown on the abstract page https://arxiv.org/abs/2506.18644v1. A full rights notice is delivered at licenses/arxiv-2506.18644-CC-BY-4.0.txt. Characters: every retained excerpt is pure ASCII, so the delivered engine, XeLaTeX with its default Latin Modern text font, reports no missing character.
  \begin{prop}
    We have \[\omega_{vect}(C_4^d)=\omega_{vect}^s(C_4^d)=2(1-1/\sqrt{3}).\]
  \end{prop}

  \begin{proof}
  The directed $4$-cycle game $(\cl G,\pi)$ has symmetries induced by $(x,y,a,b)\mapsto (x+1,y+1,a+1,b+1)$ and $(x,y,a,b)\mapsto (y,x,b,a)$. The first symmetry implies that $(x,y,a,b)\mapsto(x,y,a+1,b+1)$ and $(x,y,a,b)\mapsto (x+1,y+1,a,b)$ are also symmetries. Hence, there is an optimal density in $\cl C_{vect}$ which is invariant under these symmetries.

  For each $p\in\mathcal{C}_{vect}$, there exist vectors $v_{x,a}$,
  $w_{y,b}$ in some Hilbert space with
  $p(a,b|x,y)=\langle v_{x,a},w_{y,b}\rangle$,
  $\sum_a v_{x,a}=\eta=\sum_b w_{y,b}$ for some unit vector $\eta$,
  for all $x,y$, and the triples of vectors in each of these sums are
  mutually orthogonal. By symmetry, we may assume that, viewed as a $(\bb Z_4\times \bb Z_3)\times (\bb Z_4\times \bb Z_3)$ matrix, we have
  \[ p=\begin{pmatrix}
      X&Y&Z&Y^t\\
      Y^t&X&Y&Z\\
      Z&Y^t&X&Y\\
      Y&Z&Y^t&X
    \end{pmatrix}
  \]
  for some $3\times 3$ circulant matrices $X,Y,Z$ with entries in $[0,1]$ such that $X=X^t$ and $Z=Z^t$, and the sum of the entries of each of $X$, $Y$ and $Z$ is $1$. So $X,Y,Z\in \tfrac13\conv\{I_3,S_3,S_3^2\}$. From this block structure, we see that for $x,y\in \bb Z_4$ we have
  \begin{align*}
    X&=(\langle v_{x,a}|w_{x,b}\rangle)_{a,b},
    \\
    Y&=(\langle v_{x,a}|w_{x+1,b}\rangle)_{a,b},
    \\
    Y^t&=(\langle v_{x,a}|w_{x+3,b}\rangle)_{a,b},
    \\
    Z&=(\langle v_{x,a}|w_{x+2,b}\rangle)_{a,b}.
  \end{align*}
  Now let $v_{x,a}'=w_{x,a}'$ be given by
  \[ v_{x,a}'=w_{x,a}'=
    \begin{cases}
      v_{x,a}&\text{$x$ even,}\\
      w_{x,a}&\text{$x$ odd.}
    \end{cases}
  \]
  Let
  $p'=(\langle v'_{x,a}|v'_{y,b}\rangle)_{(x,a),(y,b)}\in \cl
  C_{vect}$ be the corresponding vector correlation. It follows that
  $p'$ is positive semidefinite and
  \[ p'
    =
    \begin{pmatrix}
      X_0&Y&Z_0&Y^t\\
      Y^t&X_1&Y&Z_1\\
      Z_0^t&Y^t&X_2&Y\\
      Y&Z_1^t&Y^t&X_3\\
    \end{pmatrix}
  \]
  for some diagonal matrices $X_x$, $x\in \bb Z_4$, and some
  $Z_0,Z_1\in M_3(\bb C)$. The point is that the matrix blocks $Y$ and $Y^t$ are
  unchanged from $p$, and the strategy value only depends on these blocks. We
  can average $p'$ over the game symmetries to get $q\in \cl C_{vect}$ with the
  same $Y$ blocks, and each diagonal block will be a diagonal
  circulant matrix with entries adding to $1$, so $\tfrac13 I_3$. Hence,
  \[ q=
    \begin{pmatrix}
      \tfrac13I_3&Y&W&Y^t\\
      Y^t&\tfrac13I_3&Y&W\\
      W&Y^t&\tfrac13I_3&Y\\
      Y&W&Y^t&\tfrac13I_3
    \end{pmatrix}\in  M_{12}(\bb C)^+\]
  for some symmetric circulant matrix $W$, and $\bb E_{\cl G,\pi}(q)=\bb E_{\cl G,\pi}(p)$. So
  $\omega_{vect}(\cl G)=\sup_q\bb E_{\cl G,\pi}(q)$ where the sup is over $q$ of this form (where $Y$ is a circulant matrix). %

  Let $u^*=2(1-\frac1{\sqrt3})\approx 0.845$. We first wish to show that $\omega_{vect}(C_4^d)\le u^*$.  Parameterize the matrices $Y$ and $W$ above as $Y=\tfrac13((1-u-v)I_3+uS_3+vS_3^2)$ for some non-negative $u,v$ with $u+v\le 1$, and $W=\tfrac13(rI_3+\frac{(1-r)}2(S_3+S_3^2))$, where $r\in [0,1]$ (recall that $W=W^t$). Note that $\bb E_{\cl G,\pi}(q)=u$. We will show that $q$ being positive semidefinite implies that $u\le u^*$.

  The $3\times 3$ blocks of $q$ are mutually commuting. The eigenvalues of $3Y$ are $\{1,\lambda,\overline{\lambda}\}$ where $\lambda=1-u-v+u\xi+v\xi^2=1-\frac32(u+v)+i\frac{\sqrt{3}}2(u-v)$, and the eigenvalues of $3W$ are $\{1,s,s\}$ where $s=(3r-1)/2$. Diagonalizing, we find that $3q\cong Q \oplus R\oplus \overline{R}$ where $Q$ is the $4\times 4$ matrix of ones, and $R$ is the $4\times 4$ circulant matrix with first row $(1,\lambda,s,\overline{\lambda})$, i.e., $R=I_4+sS_4^2+2Re(\lambda S_4)$. The eigenvalues of $R$ are
  \[
    1+(-1)^ks+2Re(i^k\lambda),\; k=0,1,2,3.
  \]
  These must all be non-negative for $q$ to be positive semidefinite.
  So $2|Re(\lambda)|\le 1+s$ and $2|Im(\lambda)|\le 1-s$.
  Recalling our paramaterization in terms of $u,v,r$, we obtain the necessary conditions
  \[ u+v\le \frac r2+\frac56,\quad u-v\le \frac{\sqrt{3}}2(1-r),\quad 0\le u,v,r\le 1.\]
  Let $r^*=\frac73-\frac4{\sqrt3}$. For each $r\in [0,1]$, the maximum value of $u$ turns out to be
  \[
    u_{max}(r)=\begin{cases}
      (5+3 r)/6&0\le r\le r^*\\
      (5+3\sqrt{3}-3(\sqrt{3}-1)r)/12&r^*\le r\le 1
    \end{cases}
  \]
  This clearly attains its maximum at $r=r^*$, so $\omega_{vect}(\mathcal{G})\le u(r^*)=u^*$.

  Finally, to see that $\omega_{vect}^s(C_4^d)\ge u^*$,
  let  $R_\theta$ be the $2\times 2$ matrix of rotation by $\theta$, for $(x,a)\in \bb Z_4\times \bb Z_3$ define unit vectors $s_{x,a},t_{x,a}\in \bb R^2$ by
   \[s_{0,0}=t_{0,0}=e_1,\quad s_{x+1,a+1}=R_{\pi/6}s_{x,a},\quad t_{x+1,a+1}=R_{\pi/3}t_{x,a},\]
   and consider the vectors $v_{x,a}\in \bb R^5$,
  \[ v_{x,a}=\frac13\begin{pmatrix}1\\ \alpha s_{x,a}\\\beta t_{x,a}
     \end{pmatrix}
   \]
   where $\alpha=\sqrt{2(\sqrt{3}-1)}$ and $\beta=\sqrt{2(2-\sqrt3)}$. Observe that
   \[ s_{x,a+1}=R_{\pi/6}^4s_{x,a}=R_{2\pi/3}s_{x,a},\quad
     t_{x,a+1}=R_{\pi/3}^4t_{x,a}=R_{4\pi/3}t_{x,a}.\] So
   \[\langle
    v_{x,a}|v_{x,a'}\rangle=\tfrac19(1+\alpha^2\cos(2\pi/3)+\beta^2\cos(4\pi/3))=\tfrac19(1-\tfrac12(\alpha^2+\beta^2))=0\]
   for $a\ne a'$.  Note that
   $\sum_{a}s_{x,a}=(I+R_{2\pi/3}+R_{4\pi/3})s_{x,0}=0$ and similarly
   $\sum_at_{x,a}=0$, so $\sum_a v_{x,a}=(1,0,0,0,0)^t$ for every
   $x$. And one can check by explicit calculation that
   $p(a,b|x,y):=\langle v_{x,a}|v_{y,b}\rangle$ has no negative
   entries. So $\{v_{x,a}\}$ defines a synchronous vect-density $p(a,b|x,y)$. Its
   winning probability for the cycle game is
   \begin{align*}\frac14&\sum_{x\in \bb Z_4,a\in \bb Z_3}p(a,a+1|x,x+1)=3p(0,1|0,1)\\&=3\langle
     v_{0,0}|v_{1,1}\rangle=\tfrac13(1+\alpha^2\cos(\pi/6)+\beta^2\cos(\pi/3))=u^*.\qedhere
   \end{align*}
 \end{proof}

\end{document}
